\section{Step 14 Insert the reconstructed color into its cache slot}\label{qoi-s14}
Author \texttt{w\_hash} to extract the previous color’s four components, form their weighted sum and select the low six bits. Convert that slot to a byte offset only when addressing its four-byte entry. Store the complete RGBA word. This common stage must run for repeated pixels too, so the first opaque-black run establishes a cache value that did not exist initially.

Work in \texttt{w\_hash}; use \texttt{s14-1} and stop at \texttt{after\_w\_hash}.

\paragraph{Observe} Cache slot fifty-three contains opaque black. The other entries remain transparent black; the caller’s metadata and framebuffer are untouched.

\paragraph{Check} Which value belongs in the cache entry?
\begin{itemize}
\item[A.] RGB with alpha forced to zero.
\item[B.] The complete RGBA pixel.
\item[C.] Only the hash number.
\end{itemize}
